\section{Résultats}
\label{sec:resultats}

Cette section présente les résultats expérimentaux de la comparaison des deux modèles sur Chuffed et Gecode.

\subsection{Performance globale}

Le Tableau~\ref{tab:resultats_temps} présente les temps de résolution pour chaque combinaison instance-modèle-solveur. Les timeouts (dépassement du temps limite de 300 secondes) sont indiqués par \textbf{TO}.

\begin{table}[h]
\centering
\caption{Temps de résolution (secondes) pour chaque combinaison instance-modèle-solveur}
\label{tab:resultats_temps}
\begin{tabular}{lcccc}
\toprule
Instance & Modèle 1 & Modèle 1 & Modèle 2 & Modèle 2 \\
         & (Chuffed) & (Gecode) & (Chuffed) & (Gecode) \\
\midrule
small\_01 & [À compléter] & [À compléter] & 0.6 & 0.3 \\
small\_02 & [À compléter] & [À compléter] & 1.9 & 4.8 \\
medium\_01 & [À compléter] & [À compléter] & \textbf{UNSAT} (1.5s) & \textbf{UNSAT} (1.4s) \\
medium\_02 & [À compléter] & [À compléter] & \textbf{UNSAT} (2.6s) & \textbf{TO} \\
medium\_03 & [À compléter] & [À compléter] & \textbf{UNSAT} (4.0s) & \textbf{TO} \\
large\_01 & [À compléter] & [À compléter] & \textbf{UNSAT} (5.9s) & \textbf{UNSAT} (26.5s) \\
large\_02 & [À compléter] & [À compléter] & \textbf{UNSAT} (14.1s) & \textbf{UNSAT} (12.8s) \\
large\_03 & [À compléter] & [À compléter] & \textbf{UNSAT} (13.5s) & \textbf{UNSAT} (12.6s) \\
\bottomrule
\end{tabular}
\end{table}

\textbf{Note :} TO = Timeout (> 300s), UNSAT = Instance insatisfiable (aucune solution n'existe). Les valeurs entre parenthèses indiquent le temps nécessaire pour prouver l'insatisfiabilité.

\subsection{Comparaison des modèles}

[À compléter avec l'analyse comparative des deux modèles : quel modèle est le plus rapide en moyenne, y a-t-il des instances où un modèle excelle et l'autre échoue, quel est l'impact de la modélisation sur la performance ?]

\subsection{Comparaison des solveurs}

[À compléter avec l'analyse comparative de Chuffed vs Gecode : pour chaque modèle, quel solveur est le plus performant, y a-t-il des différences selon le modèle utilisé, quel est le taux de succès par solveur et par modèle ?]

\subsection{Observations principales}

[À compléter avec les observations clés : à partir de quelle taille les solveurs commencent-ils à échouer, y a-t-il des patterns dans les instances difficiles, quel est l'impact de la modélisation sur la résolubilité ?]

